%%═══════════════════════════════════════════════════════════════════
%% Lecture 10 — Multidimensional Root Finding
%% 77315 — Computational Physics
%%═══════════════════════════════════════════════════════════════════

\section{Lecture 10 — Multidimensional Root Finding}\label{sec:lec10}
\hrule\vspace{1.5em}

%%──────────────────────────────────────────────────────────────────
\subsection*{Overview}
Many physical equilibrium problems require solving $n$ simultaneous
nonlinear equations $\mathbf{F}(\mathbf{x})=\mathbf{0}$ for
$\mathbf{x}\in\mathbb{R}^n$.  The 1D Newton–Raphson update generalises
directly to $n$ dimensions using the Jacobian matrix.  However, the
Newton step can overshoot; we introduce a \emph{merit function} and
\emph{line-search backtracking} to guarantee descent.

%%──────────────────────────────────────────────────────────────────
\subsection{Multidimensional Newton--Raphson}\label{subsec:lec10:newtonraphson}

\subsubsection*{Setup}
Given $n$ equations $\mathbf{F}(\mathbf{x})=\mathbf{0}$, define the
Jacobian matrix
\begin{equation}
  J_{ij}(\mathbf{x}) = \frac{\partial F_i}{\partial x_j}.
  \label{eq:lec10:jacobian}
\end{equation}
Linearising $\mathbf{F}$ around $\mathbf{x}_k$:
\begin{equation}
  \mathbf{F}(\mathbf{x}_k + \boldsymbol{\delta}) \approx
  \mathbf{F}(\mathbf{x}_k) + J(\mathbf{x}_k)\,\boldsymbol{\delta} = \mathbf{0}
  \;\Rightarrow\;
  J\,\boldsymbol{\delta} = -\mathbf{F}.
  \label{eq:lec10:linear}
\end{equation}

\subsubsection*{Newton--Raphson update}
\begin{equation}
  \boxed{\mathbf{x}_{k+1} = \mathbf{x}_k - J^{-1}(\mathbf{x}_k)\,\mathbf{F}(\mathbf{x}_k)}
  \label{eq:lec10:NR}
\end{equation}
In practice we never invert $J$; instead we solve the linear system
$J\,\boldsymbol{\delta} = -\mathbf{F}$ (e.g.\ via LU decomposition)
and then set $\mathbf{x}_{k+1} = \mathbf{x}_k + \boldsymbol{\delta}$.

The method converges \emph{quadratically} near the root:
$\|\mathbf{x}_{k+1}-\mathbf{x}^*\| = \bigO{\|\mathbf{x}_k - \mathbf{x}^*\|^2}$.

\subsubsection*{Algorithm}
\begin{enumerate}
  \item Choose initial guess $\mathbf{x}_0$ and tolerance $\eps$.
  \item Evaluate $\mathbf{F}(\mathbf{x}_k)$; if $\|\mathbf{F}\|<\eps$ stop.
  \item Compute $J(\mathbf{x}_k)$ (analytically or by finite differences).
  \item Solve $J\,\boldsymbol{\delta} = -\mathbf{F}$ for $\boldsymbol{\delta}$.
  \item Set $\mathbf{x}_{k+1} = \mathbf{x}_k + \boldsymbol{\delta}$ and go to 2.
\end{enumerate}

\nt{A poor initial guess can cause divergence or convergence to an unwanted root.  Scaling variables to $\bigO{1}$ improves robustness.}

%%──────────────────────────────────────────────────────────────────
\subsection{Merit Function and Line-Search Backtracking}\label{subsec:lec10:linesearch}

\subsubsection*{Merit function}
The Newton step can overshoot.  Introduce the scalar merit function
\begin{equation}
  f(\mathbf{x}) = \tfrac{1}{2}\,\mathbf{F}(\mathbf{x})\cdot\mathbf{F}(\mathbf{x})
                = \tfrac{1}{2}\|\mathbf{F}\|^2.
  \label{eq:lec10:merit}
\end{equation}
A successful iteration must \emph{decrease} $f$, i.e.\ it must move
downhill on the merit-function landscape.

\subsubsection*{Line search}
Instead of accepting the full Newton step $\boldsymbol{\delta}$,
move along the Newton direction $\hat{\boldsymbol{\delta}}$ with a
step-length $\lambda\in(0,1]$:
\begin{equation}
  \mathbf{x}_{\text{new}} = \mathbf{x}_k + \lambda\,\boldsymbol{\delta}.
  \label{eq:lec10:step}
\end{equation}
Start with $\lambda=1$ (full Newton step).  If
$f(\mathbf{x}_{\text{new}}) > f(\mathbf{x}_k)$, \emph{backtrack}:
reduce $\lambda$ (e.g.\ $\lambda \leftarrow \lambda/2$) and try again.
Repeat until $f$ decreases or $\lambda$ becomes negligibly small.

\nt{Backtracking ensures global convergence from far-from-root starting points where pure Newton steps may not decrease $\|\mathbf{F}\|$.}

%%──────────────────────────────────────────────────────────────────
\subsection{Jacobian by Finite Differences}\label{subsec:lec10:fd_jacobian}

When the analytic Jacobian is unavailable, approximate each column
by a forward difference:
\begin{equation}
  J_{ij} \approx \frac{F_i(\mathbf{x}+h\,\hat{e}_j)-F_i(\mathbf{x})}{h},
  \label{eq:lec10:fdJ}
\end{equation}
at the cost of $n$ extra function evaluations per Newton step.
The step $h$ must balance truncation and round-off; a typical choice
is $h = \eps_\text{mach}^{1/2}\max(1,|x_j|)$.
